% GENERATED FILE. Edit canonical lessons, then run npm run generate:books.
% canonical-source-hash: fnv1a64:fb3379123cca4e37
% canonical-lessons: HI-C188-batan, HI-C188-jeb, HI-C188-dusman, HI-C188-caci, HI-C188-bhatija

\chapter{A button, A pocket, An enemy, An aunt (father's brother's wife), A nephew (brother's son)}
\label{ch:hi-a-button-a-pocket-an-enemy-an-aunt-father-s-brother-s-wife-a-nephew-brother-s-son}
\hlchaptermodality{188}

\begin{chapteropening}
\textbf{By the end of this chapter:} \emph{I can say a button, a pocket, an enemy, an aunt (father's brother's wife) and a nephew (brother's son).}

Complete the last lesson of chapter 188: i can say a button, a pocket, an enemy, an aunt (father's brother's wife) and a nephew (brother's son).
\end{chapteropening}

\section[\texorpdfstring{\hi{बटन} (baṭan)}{बटन (baṭan)}]{\hi{बटन} (baṭan) — a button}

\label{lesson:HI-C188-batan}



\begin{quote}
Before the new one: say the Hindi for the cinema, then the Hindi for a park.
\end{quote}

\subsection*{You'll want to know: \hi{बटन}}
\textbf{\hi{बटन}} — \emph{baṭan} — \textquotedblleft{}a button\textquotedblright{}.

\begin{grammarlens}[title={What you've built}]
One more word for this chapter.
\end{grammarlens}

\subsection*{Your turn}
\begin{itemize}
\raggedright
  \item \emph{Say it:} \emph{baṭan}
  \item \emph{Say it:} \emph{baṭan}, once more
  \item \emph{Say it:} say \emph{baṭan}
  \item \emph{Recall:} say the Hindi for the cinema, then the Hindi for a park, then say \emph{baṭan} again
\end{itemize}

\begin{tcolorbox}[breakable,colback=teal!4,colframe=teal!35!black,title={Before you move on}]
What does \hi{बटन} mean? (\textquotedblleft{}A button\textquotedblright{}.) Say it once more.
\end{tcolorbox}
\section[\texorpdfstring{\hi{जेब} (jeb)}{जेब (jeb)}]{\hi{जेब} (jeb) — a pocket}

\label{lesson:HI-C188-jeb}



\begin{quote}
Before the new one: say the Hindi for a park, then the Hindi for a button.
\end{quote}

\subsection*{You'll want to know: \hi{जेब}}
\textbf{\hi{जेब}} — \emph{jeb} — \textquotedblleft{}a pocket\textquotedblright{}.

\begin{grammarlens}[title={What you've built}]
One more word for this chapter.
\end{grammarlens}

\subsection*{Your turn}
\begin{itemize}
\raggedright
  \item \emph{Say it:} \emph{jeb}
  \item \emph{Say it:} \emph{jeb}, once more
  \item \emph{Say it:} say \emph{jeb}
  \item \emph{Recall:} say the Hindi for a park, then the Hindi for a button, then say \emph{jeb} again
\end{itemize}

\begin{tcolorbox}[breakable,colback=teal!4,colframe=teal!35!black,title={Before you move on}]
What does \hi{जेब} mean? (\textquotedblleft{}A pocket\textquotedblright{}.) Say it once more.
\end{tcolorbox}
\section[\texorpdfstring{\hi{दुश्मन} (duśman)}{दुश्मन (duśman)}]{\hi{दुश्मन} (duśman) — an enemy}

\label{lesson:HI-C188-dusman}



\begin{quote}
Before the new one: say the Hindi for a button, then the Hindi for a pocket.
\end{quote}

\subsection*{You'll want to know: \hi{दुश्मन}}
\textbf{\hi{दुश्मन}} — \emph{duśman} — \textquotedblleft{}an enemy\textquotedblright{}.

\begin{grammarlens}[title={What you've built}]
One more word for this chapter.
\end{grammarlens}

\subsection*{Your turn}
\begin{itemize}
\raggedright
  \item \emph{Say it:} \emph{duśman}
  \item \emph{Say it:} \emph{duśman}, once more
  \item \emph{Say it:} say \emph{duśman}
  \item \emph{Recall:} say the Hindi for a button, then the Hindi for a pocket, then say \emph{duśman} again
\end{itemize}

\begin{tcolorbox}[breakable,colback=teal!4,colframe=teal!35!black,title={Before you move on}]
What does \hi{दुश्मन} mean? (\textquotedblleft{}An enemy\textquotedblright{}.) Say it once more.
\end{tcolorbox}
\section[\texorpdfstring{\hi{चाची} (cācī)}{चाची (cācī)}]{\hi{चाची} (cācī) — an aunt (father's brother's wife)}

\label{lesson:HI-C188-caci}



\begin{quote}
Before the new one: say the Hindi for a pocket, then the Hindi for an enemy.
\end{quote}

\subsection*{You'll want to know: \hi{चाची}}
\textbf{\hi{चाची}} — \emph{cācī} — \textquotedblleft{}an aunt (father's brother's wife)\textquotedblright{}.

\begin{grammarlens}[title={What you've built}]
One more word for this chapter.
\end{grammarlens}

\subsection*{Your turn}
\begin{itemize}
\raggedright
  \item \emph{Say it:} \emph{cācī}
  \item \emph{Say it:} \emph{cācī}, once more
  \item \emph{Say it:} say \emph{cācī}
  \item \emph{Recall:} say the Hindi for a pocket, then the Hindi for an enemy, then say \emph{cācī} again
\end{itemize}

\begin{tcolorbox}[breakable,colback=teal!4,colframe=teal!35!black,title={Before you move on}]
What does \hi{चाची} mean? (\textquotedblleft{}An aunt (father's brother's wife)\textquotedblright{}.) Say it once more.
\end{tcolorbox}
\section[\texorpdfstring{\hi{भतीजा} (bhatījā)}{भतीजा (bhatījā)}]{\hi{भतीजा} (bhatījā) — a nephew (brother's son)}

\label{lesson:HI-C188-bhatija}



\begin{quote}
Before the new one: say the Hindi for an enemy, then the Hindi for an aunt.
\end{quote}

\subsection*{You'll want to know: \hi{भतीजा}}
\textbf{\hi{भतीजा}} — \emph{bhatījā} — \textquotedblleft{}a nephew (brother's son)\textquotedblright{}.

\begin{grammarlens}[title={What you've built}]
One more word for this chapter.
\end{grammarlens}

\subsection*{Your turn}
\begin{itemize}
\raggedright
  \item \emph{Say it:} \emph{bhatījā}
  \item \emph{Say it:} \emph{bhatījā}, once more
  \item \emph{Say it:} say \emph{bhatījā}
  \item \emph{Recall:} say the Hindi for an enemy, then the Hindi for an aunt, then say \emph{bhatījā} again
\end{itemize}

\begin{tcolorbox}[breakable,colback=teal!4,colframe=teal!35!black,title={Before you move on}]
What does \hi{भतीजा} mean? (\textquotedblleft{}A nephew (brother's son)\textquotedblright{}.) Say it once more.
\end{tcolorbox}
